%==================================================================
% ABSTRACT (Bahasa Inggris)
% Acuan: Pedoman Lampiran 9B
%==================================================================
\clearpage
\thispagestyle{plain}
\frontmatterentry{ABSTRACT}

\begin{spacing}{1}
\begin{center}
{\fontsize{14}{16.8}\selectfont\bfseries\itshape
 \MakeUppercase{\judulen}\par}
\end{center}

\vspace{2\baselineskip}

\noindent
\begin{tabular}{@{}l@{\hspace{0.6em}}c@{\hspace{0.6em}}p{9.2cm}@{}}
By                          & : & \penulis \\
Student Identification Number & : & \nim \\
Supervisor                  & : & 1.\ \pembimbingA \\
                            &   & 2.\ \pembimbingB \\
\end{tabular}

\vspace{2\baselineskip}

\begin{center}
{\fontsize{14}{16.8}\selectfont\bfseries\itshape ABSTRACT\par}
\end{center}

\vspace{2\baselineskip}

%--- Isi abstrak: terjemahan dari abstrak Bahasa Indonesia ----------
The decline of fish freshness occurs very quickly after catching, so a
fast and objective freshness assessment is essential for both buyers
and fishery business actors. At present, freshness is still assessed
organoleptically by trained personnel, which makes the result
subjective and difficult to perform at scale. This study aims to design
a fish freshness detection system based on digital image processing and
an artificial neural network that is able to classify fish freshness
automatically. Images of fish eyes and gills were acquired using a
digital camera inside a capture box with fixed illumination, then
processed through resizing, contrast enhancement, and colour feature
extraction in the HSV colour space. The extracted features were used as
inputs to a backpropagation neural network with two hidden layers.
Testing was carried out on 300 fish images covering three freshness
classes, namely fresh, moderately fresh, and not fresh. The test result
shows a classification accuracy of 93.33\% with an average processing
time of 0.42 seconds per image. The highest precision and recall were
obtained for the fresh class, while most errors occurred between the
moderately fresh and not fresh classes due to similar gill colours.
Therefore, the designed system can be used as an alternative for fast,
objective, and non-destructive fish freshness assessment.

\vspace{3\baselineskip}

\noindent
\textbf{Keywords:} artificial neural network, digital image processing,
fish freshness, HSV colour space
%-------------------------------------------------------------------
\end{spacing}
